Esta monografia especifica, implementa e verifica o subsistema de memória cache de nível 1 de um processador RISC-V descrito em Chisel, em um núcleo que visa executar o Linux. O método compreendeu a especificação do contrato do subsistema em sinais e temporização, independente da linguagem; a implementação das caches de instrução e de dados como instâncias parametrizadas de um mesmo gerador; e a verificação por 180 casos automatizados, comparados em tempo de execução contra um modelo de referência escrito independentemente em Scala. A avaliação submeteu treze configurações de geometria a um trace fixo de 1.500 acessos, variando capacidade, associatividade, tamanho de linha e política, e instrumentou o custo da instrução de sincronização. Os resultados mostram taxa de acerto de 69,5\% na configuração adotada, de 8 KiB com quatro vias e linha de 16 bytes, que satura nesse ponto: dobrar a capacidade acrescenta um único acerto em 1.500. Linhas de 64 bytes acertam 14,7 pontos mais, mas transferem 2,5 vezes mais bytes. Conclui-se que a configuração adotada minimiza o tráfego dentro da faixa de acerto aceitável, e que a taxa de acerto isolada é critério insuficiente quando o custo da transferência acompanha o volume.
